% Generated draft: requires mathematical and source review.
\documentclass[11pt]{article}
\usepackage[margin=25mm]{geometry}
\usepackage{fontspec}
\setmainfont{Latin Modern Roman}
\usepackage{amsmath,amssymb,mathtools,unicode-math}
\setmathfont{Latin Modern Math}
\usepackage{enumitem,needspace,xurl,hyperref}
\setlength{\parindent}{0pt}
\setlength{\parskip}{5pt}
\setlength{\emergencystretch}{3em}
\newcommand{\N}{\mathbb{N}}
\newcommand{\Z}{\mathbb{Z}}
\newcommand{\Q}{\mathbb{Q}}
\newcommand{\R}{\mathbb{R}}
\newcommand{\C}{\mathbb{C}}
\newcommand{\F}{\mathbb{F}}
\newcommand{\A}{\mathbb{A}}
\newcommand{\PP}{\mathbb{P}}
\newcommand{\E}{\mathbb{E}}
\newcommand{\eps}{\varepsilon}
\newcommand{\dd}{\,\mathrm{d}}
\begin{document}
\section*{UnsolvedMath ID: 20001643}
\textbf{Source problem code:} catalogue entry
\subsection{Definitions and mathematical statement}
Let \(F_n\) be free of rank \(n\) and \(\operatorname{Out}(F_n)=\operatorname{Aut}(F_n)/\operatorname{Inn}(F_n)\). Property \((T)\) means that every unitary representation with almost-invariant unit vectors has a nonzero invariant vector. For a fixed finite symmetric generating multiset \(S\), Property \((\tau)\) relative to a family \(\mathcal F\) of finite-index subgroups means a uniform positive spectral gap on
\(\ell^2(H\backslash G)\ominus\C\mathbf1\), \(H\in\mathcal F\), for the normalized Schreier Laplacian. The source asks whether \(\operatorname{Out}(F_n)\) has \((T)\) or \((\tau)\) for \(n\ge4\). A family must be specified for \((\tau)\); ``all finite-index subgroups'' is the strong interpretation used in the supplied aid.

The historical motivation states that \(\operatorname{Out}(F_3)\) has a finite-index subgroup with infinite abelianization and hence lacks \((T)\). The aid reports an affirmative answer for all \(n\ge4\): results for \(\operatorname{Aut}(F_4)\), \(\operatorname{Aut}(F_5)\), and \(\operatorname{Aut}(F_n)\), \(n\ge6\), pass to the outer quotient, and \((T)\) implies \((\tau)\) for every finite-index family. It expressly distinguishes a computer-assisted rank-four arXiv preprint, whose attached Sage verifier was not rerun, from peer-reviewed rank-five and higher-rank inputs. No identifiers or complete references for these inputs are supplied, so this is a qualified literature report, not an independently verified certificate.

The aid proves a quantitative quotient transfer. For a finitely generated group \(G\), put
\[
 \Delta_S=|S|\,1-\sum_{s\in S}s\in\R[G],
\]
and use the involution \(g^*=g^{-1}\). If an exact identity
\[
 \Delta_S^2-\epsilon\Delta_S=\sum_j\xi_j^*\xi_j,\quad \epsilon>0,
\]
holds, then a quotient \(q:G\twoheadrightarrow Q\) gives the same identity with parameter \(\epsilon\) for the projected multiset \(q(S)\) and \(q(\xi_j)\). Multiplicities and identity images are retained. Thus the unnormalized gap parameter is unchanged; deleting repetitions or changing Laplacian normalization is a separate operation.

Conversely, suppose a finite-index subgroup \(H\le G\) has an epimorphism \(\chi:H\to\Z\). For integers \(q\ge2\), put \(H_q=\chi^{-1}(q\Z)\). Choose representatives \(r\) for \(H\backslash G\), write \(rs=c(r,s)r'\) with \(c(r,s)\in H\), and put \(M=\max_{r,s}|\chi(c(r,s))|\). The reported normalized Schreier-gap bound is
\[
 \lambda_1(H_q\backslash G)\le\frac{2\pi^2M^2}{q^2}.
\]
The constant is independent of \(q\). Pullback to the normal cores \(C_q=\bigcap_{g\in G}gH_qg^{-1}\) gives the same failure of uniform gap for an all-normal finite-index family; choosing divisibility-nested \(q\)'s makes a tower. No trivial-intersection assertion is needed. In particular virtual infinite abelianization obstructs \((\tau)\) for the all-finite-index family as well as \((T)\), but says nothing by itself about an arbitrarily restricted family.
\subsection{Short English statement}
The supplied literature report answers the historical rank-at-least-four question affirmatively, while flagging the unrerun computer-assisted rank-four preprint. Its quantitative lemmas preserve exact sum-of-squares gaps through quotients and turn virtual infinite abelianization into vanishing finite-quotient Schreier gaps.
\subsection{Sources}
\begin{itemize}
\item AIM workshop \emph{The geometry of the outer automorphism group of a free group}, Problem 3.1 and its rank-three motivation: \url{https://aimath.org/WWN/outerauto/outerauto.pdf}.
\item Local registry \texttt{data/unsolvedmath/problems.json}, permanent ID 20001643, source code \texttt{catalogue entry}; canonical source \texttt{aim-geometric-group-theory-notes.json}, \texttt{notes[151]}.
\item The supplied summary cites unnamed Property \((T)\) inputs for \(\operatorname{Aut}(F_4)\), \(\operatorname{Aut}(F_5)\) and higher ranks: the first is a computer-assisted arXiv preprint with a Sage verifier not rerun in that attempt; the others are reported peer reviewed. Bibliographic identifiers are not supplied.
\item Attached machine-generated attempt 1: full solution, run status solved, internal rigor label ``proved'', low novelty confidence; not peer reviewed and requiring independent verification before citation.
\end{itemize}
\end{document}
